\documentclass[12pt]{article}
\usepackage{amsfonts,amsthm,amsmath,amssymb,mathtools}
\usepackage{mathrsfs}
\usepackage{enumitem}
\usepackage{tikz-cd}
\usepackage{geometry}
\usepackage{mlmodern}
\usepackage{graphicx}

\geometry{letterpaper, portrait, margin=1in}

\setcounter{section}{0}

\renewcommand{\thesection}{\S\arabic{section}}
\renewcommand{\thesubsection}{\arabic{section}}
\DeclareMathOperator{\im}{im}
\DeclareMathOperator{\id}{id}

\newtheorem{thm}{Theorem}
\newenvironment{theorem}[1]{
	\renewcommand{\thethm}{#1}
	\begin{thm}
		}{\end{thm}}

\newtheorem{lma}{Lemma}
\newenvironment{lemma}[1]{
	\renewcommand{\thelma}{#1}
	\begin{lma}
		}{\end{lma}}

\theoremstyle{definition}
\newtheorem{ex}{Assignment}
\newenvironment{exercise}[1]{
	\renewcommand{\theex}{#1}
	\begin{ex}
		}{\end{ex}}

\title{MTH 732 Topology - Homework 5}
\author{Connor Mooneyhan}
\date{February 19, 2026}

\begin{document}

\maketitle

\begin{ex}
	Find all connected covering spaces of $\mathbb{RP}^2 \vee \mathbb{RP}^2$.
\end{ex}
\begin{proof}
	First, we have the universal cover, which is an infinite strip of spheres.
  For instance, you could model in $\mathbb{R^3}$ with spheres of radius $0.5$, each centered at $(n, 0, 0)$ where $n$ is an integer, so that there is a wedge point at each $(x, 0, 0)$ where $x - 0.5 \in \mathbb{Z}$.
  Call this universal cover $E$.
  Every other connected covering space will be a chain of $n$ spheres, formed by the quotient $E/((x, 0, 0) \sim (x + nz, 0, 0))$.
\end{proof}

\begin{ex}
	Find a space $E$ for which there exists a 3-sheeted covering map $p: E \to T^2$ of the torus $T^2$.
	Describe the mapping $p$.
\end{ex}
\begin{proof}
	Let $E$ be $[0, 1] \times [0, 3]$, where $(x, 0) \sim (x, 3) \forall x \in [0, 1]$ and $(0, y) \sim (1, y) \forall y \in [0, 3]$.
	We will realize $T^2$ as the usual identification space on $\mathbb{R}^2/\mathbb{Z}^2$.
	The covering map $p : E \to T^2$ is then described as \begin{equation*}
		p([(x, y)]_E) = ([x, y]_{T^2}).
	\end{equation*}
	We see that for any point $([x, y]_{T^2}) \in \mathbb{R}^2/\mathbb{Z}^2$, we have \begin{equation*}
		p^{-1}([(x, y)]_{T^2}) = \lbrace[(x, y)], [(x, y + 1)], [(x, y + 2)]\rbrace.
	\end{equation*}
	For any $[(x, y)]_{T_2} \in T^2$, then, there is a neighborhood $U$ such that $p^{-1}(U)$ is the union of three disjoint open neighborhoods of elements of $p^{-1}([(x, y)]_{T^2})$, each of which is homeomorphic to $U$.
	Thus, $p$ is a covering map.
\end{proof}

\begin{ex}
	Let $X$ and $Y$ be path-connected, locally path-connected, and semilocally simply connected spaces.
	Let $\widetilde{X}$ and $\widetilde{Y}$ be simply connected covers of  $X$ and $Y$, respectively.
	If $X$ is homotopy equivalent to $Y$, show that $\widetilde{X}$ is homotopy equivalent to $\widetilde{Y}$.
\end{ex}
\begin{proof}
  Let $f : X \to Y$ be a homotopy equivalence and let $g : Y \to X$ be a homotopy inverse of $f$.
  Also, let $p_X : \widetilde{X} \to X$ and $p_Y : \widetilde{Y} \to Y$ be the covering maps.
  Then $f \circ p_X : \widetilde{X} \to Y$ lifts (since $\pi_1(\widetilde{X})$ is trivial and thus its image under $(f \circ p_X)_*$ is trivial) to a map $\widetilde{f} : \widetilde{X} \to \widetilde{Y}$.
  Similarly, $g \circ p_Y$ lifts to a map $\widetilde{g} : \widetilde{Y} \to \widetilde{X}$.

  Now since $f \circ g \simeq \id_Y$, $f \circ g \circ p_Y \simeq p_Y$.
  This homotopy can then be lifted along $p_Y$ to get $\widetilde{f} \circ \widetilde{g} \sim \id_{\widetilde{Y}}$
  Similarly, we get $\widetilde{g} \circ \widetilde{f} \sim \id_{\widetilde{X}}$.
  Thus $\widetilde{X}$ and $\widetilde{Y}$ are homotopy equivalent.
\end{proof}

\begin{ex}
	Show that a torus cannot cover a sphere, and vice-versa.
\end{ex}
\begin{proof}
	Let $p: T^2 \to S^2$ be a covering map.
	Then $p$ induces an injection from $\pi_1(T^2)$ to $\pi_1(S^2)$, which is impossible since $\pi_1(T^2)$ is non-trivial and $\pi_1(S^2)$ is trivial.
  In the other direction, if $S^2$ covers $T^2$, then since $S^2$ is simply connected, it would be a universal cover, and thus it would be homeomorphic to the already-known universal cover $\mathbb{R}^2$.
  Since $S^2$ is not homeomorphic to $\mathbb{R}^2$, we are done.
\end{proof}

\begin{ex}
	Describe a covering map $p : M_3 \to M_2$, where $M_g$ is the connected, compact surface of genus $g$.
\end{ex}
\begin{proof}
  Not sure what to do here.
\end{proof}

\end{document}
